% =====================================================================
% Mecanum_PINN_Motivation_RelatedWork.tex
%
% Motivation + Related-Work section for the IMECE 2026 paper on a
% PINN + Mamba digital twin for forward-dynamics learning and per-wheel
% friction-parameter (mu, chi) recovery on a KUKA youBot four-mecanum
% platform.
%
% This file compiles STANDALONE (pdflatex + bibtex) for review, and the
% two \section blocks lift directly into the ASME imece template.
%
% To drop into the ASME template:
%   1. Delete the preamble down to "% ===== SECTION CONTENT BEGINS".
%   2. Paste the two sections after your \section{Introduction}.
%   3. Use the ASME .bst (e.g. asmems4) instead of unsrtnat below, and
%      \bibliography{mecanum_pinn_refs}.
%
% EVIDENCE NOTE: every quantitative figure below passed 3-vote adversarial
% verification. Two standing caveats, stated once here and reflected in the
% wording: (a) all reported gains are the cited authors' own results against
% their own baselines, not independently replicated; (b) most PINN
% friction-ID evidence comes from FIXED-BASE manipulators, not mecanum
% mobile robots -- their transfer to per-wheel (mu, chi) recovery is by
% analogy, which is precisely the gap this paper closes.
% =====================================================================
\documentclass[11pt]{article}

\usepackage[utf8]{inputenc}
\usepackage[T1]{fontenc}
\usepackage{amsmath,amssymb}
\usepackage[margin=1in]{geometry}
\usepackage{booktabs}
\usepackage{array}
\usepackage[numbers,sort&compress]{natbib}
\usepackage{hyperref}
\usepackage{xcolor}
\hypersetup{colorlinks=true, linkcolor=blue!50!black, citecolor=blue!50!black,
            urlcolor=blue!50!black}

\title{\bfseries Motivation and Related Work (draft)\\
\normalsize PINN + Mamba Digital Twin for Mecanum Forward Dynamics and Friction Identification}
\author{\normalsize IMECE 2026 --- section draft}
\date{\today}

\begin{document}
\maketitle

% ===== SECTION CONTENT BEGINS =========================================

\section{Motivation}
\label{sec:motivation}

Four-wheel Mecanum platforms are attractive across industrial and service
settings---warehouse and intralogistics transport, manufacturing and assembly
feeding, hospital and clinical logistics, cleaning, inspection, and assistive
mobility---because their omnidirectional maneuverability lets them translate and
rotate independently in cluttered, human-shared, space-constrained environments
where conventional differential-drive vehicles require multi-point turns
\cite{burkacki2025review}. Yet a recent systematic review of Mecanum and omni-wheel
technologies identifies four recurring barriers that gate real-world deployment:
limited stability and maneuverability on uneven and inclined terrain, high energy
consumption, complex control requirements, and elevated cost \cite{burkacki2025review}.

\subsection{Why higher-fidelity control is needed}
\label{sec:why-fidelity}

Underneath those application-level barriers sits a common set of physics-level
obstacles that the ideal-rolling kinematic model omits, and each has been shown to
cause a concrete, independently reported failure. \emph{Energy:} in a real
industrial haul---five tonnes of zinc ingots transported 5~km in an iron--steel
plant---wheel slippage caused $\approx$6.18\% energy loss, cutting battery-supported
range proportionally, which a kinematic controller assuming ideal rolling cannot
account for \cite{bayar2025battery}. \emph{Accuracy:} roller--surface contact
produces significant deviations in vehicle orientation during tracking, requiring a
control-layer intervention (a ramp-shaped reference) beyond baseline tracking to
suppress \cite{jmst2024slippage}; and the resulting odometric drift is
surface- and velocity-dependent, so the same kinematic controller yields different
positioning error across floors and speeds \cite{shahidi2021odometry}.
\emph{Coupled accuracy/energy:} restricted force/torque authority causes transient
overshoot that simultaneously lowers tracking performance and raises energy use,
reducing autonomy \cite{antipeaking2024}. \emph{Root cause:} Mecanum friction is
comparatively undermodeled relative to pneumatic or railroad wheels, and an
experimentally validated orthotropic-friction model shows the contact behavior
depends on both velocity and the payload-shifted center of mass
\cite{manzl2023orthotropic}. Because these effects are nonlinear, load-dependent,
and surface-dependent, purely kinematic or partially dynamic controllers cannot
reject them---precisely why higher-fidelity model-based control and system
identification are required \cite{li2025mecanumdob,avpsmo2024}.

Most directly for this work, on the \emph{KUKA youBot} platform itself the real
Mecanum-wheel construction (a finite number of rollers and their arrangement)
degrades odometric navigation whenever the odometry algorithm treats the wheel as
an idealized disc, and an algorithm that accounts for the true geometry measurably
reduces the coordinate error \cite{adamov2018youbot}. This on-platform evidence
motivates the friction/dynamics fidelity pursued here (the simulator's LuGre--Adamov
friction model follows the same line of youBot-specific modeling).

\subsection{Sector barriers and where control is decisive}
\label{sec:sector-barriers}

The same physics surfaces as distinct barriers across deployment sectors.
In \emph{heavy industry}, slip-driven energy loss shrinks effective range on
heavy-load hauls \cite{bayar2025battery}. In \emph{automotive} intralogistics,
baseline controllers lack the robustness for a four-Mecanum autonomous vehicle,
motivating adaptive sliding-mode control with an extended state observer
\cite{bui2025hybrid}. In \emph{warehouse} logistics, surface/velocity-dependent
odometry drift undermines the pose repeatability that automated workflows depend on
\cite{shahidi2021odometry}. In \emph{aerospace} assembly, tolerances on the order
of $0.01$~mm coexist with mobile robots whose position-control accuracy and
structural stiffness are low, so deployment remains largely at the proof-of-concept
stage and relies on two-stage positioning---coarse by the platform wheels, fine by a
laser-tracker-guided arm \cite{aerospaceassembly2026}. In \emph{hospitals}, mobile
service robots must satisfy six risk dimensions spanning environmental complexity,
hygiene, human interaction, workflow flexibility, and autonomy
\cite{hospitalrisks2024}, and their navigation testing is still preliminary and not
standardized \cite{hospitalnavbench2024}.

A recurring pattern ties these together. The commercial precision bar is set by
platforms such as KUKA omniMove, which moves loads up to 90~tonnes and states
contactless positioning to $\pm3$~mm \cite{kuka_omnimove}; yet its base drive
accuracy is on the order of $\pm5$~mm and is tightened to $\approx\pm1$~mm only with
add-on optical tracking, positioning aids, and mechanical guides
\cite{kuka_aerospace}---the same external-metrology workaround seen in aerospace,
where docking accuracy and repeatability are treated as a first-class requirement
\cite{nist2016docking}. In other words, deployed precision is bought by external
referencing rather than by higher-fidelity onboard control, and Mecanum
slip/odometry drift is exactly what forces that dependence
\cite{shahidi2021odometry,adamov2018youbot}. The barriers thus split cleanly:
\emph{control-related}---slip-induced orientation error, odometric drift,
disturbance/uncertainty rejection, and docking repeatability---versus
\emph{non-control}---cost, battery life, systems integration, standards compliance,
hygiene, and workforce. This paper targets the former.

\subsection{Deployability near humans favors interpretable, verifiable models}
\label{sec:certifiability}

A further, deployment-blocking barrier is regulatory rather than merely
performance-related. Driverless industrial trucks---the category covering Mecanum
AGVs and AMRs---are governed by ISO~3691-4:2023, which mandates safety functions
rated to Performance Levels (typically PL~d, corresponding to SIL~2), personnel
detection, and speed-limited operating zones; it inherits its method from
ISO~13849-1 and IEC~61508, and human-proximity collaboration adds ISO~10218 and
ISO/TS~15066 \cite{iso3691_4,iso13849}. Against this backdrop, systematic reviews
find that neural-network (black-box) controllers are confined to the lowest safety-
integrity levels precisely because their outputs cannot be exhaustively predicted or
formally verified, with interpretability a named open problem
\cite{zhang2020nnverif,certifyml2022}. This directly favors a \emph{physics-informed},
structurally-constrained model---one that retains a known friction structure and
identifies interpretable per-wheel parameters---over an unconstrained black box, both
because it is more amenable to the verification and PL rating these standards demand
and because the resulting twin can itself serve as a verification and reachability
tool for the deployed controller.

The core difficulty underlying all of the above is friction. Slip and stick--slip
behavior at the roller contact are governed by dynamic friction that simplified
constant- or Coulomb-plus-viscous models do not capture, and whose parameters are
neither directly measurable nor constant across surfaces and payloads. A digital
twin that both predicts forward dynamics at high fidelity and \emph{recovers} the
underlying per-wheel friction parameters therefore addresses the root cause rather
than compensating for its symptoms downstream---the objective of this work.

\section{Related Work}
\label{sec:related}

\subsection{Higher-fidelity model-based control closes measurable gaps}
\label{sec:related-control}

A consistent body of recent Mecanum-specific results shows that improving model
fidelity and adding explicit disturbance estimation measurably improves tracking
and robustness (Table~\ref{tab:control}). A complete Lagrangian model that
linearizes the PWM-induced motor delay reduces open-loop Dynamic-Time-Warping
(DTW) trajectory error to $0.2662$~m versus $0.3198$~m for a traditional model
($\sim$17\% improvement) \cite{li2025mecanumdob}. Coupling Internal Model Control
with a Disturbance Observer (IMC+DOB) stabilizes 97.6\% faster than IMC and 98.3\%
faster than PID under load and noise disturbance \cite{li2025mecanumdob}. Active
Disturbance Rejection Control (ADRC), which estimates and compensates unknown
dynamics with minimal plant knowledge, improves inclined-surface tracking accuracy
by 49--60\% over a PD baseline with experimental validation
\cite{hernandez2025inclined}. Adaptive and observer-based schemes go further by
estimating the disturbing quantities themselves: a Model Reference Adaptive Control
law with Lyapunov-based adaptation compensates unknown CoG offset and slope in real
time \cite{chaichumporn2025adaptive}, and an adaptive variable-power sliding-mode
observer used as a feedforward compensator within an MPC (AVPSMO-MPC) suppresses
external disturbances and model uncertainty with a Lyapunov-proven stability
guarantee \cite{avpsmo2024}. The common thread is that the residual physics---
slip, friction, and load---must be modeled or estimated, not ignored; all reported
gains are the respective authors' own results against their own baselines.

\begin{table}[t]
  \centering
  \small
  \renewcommand{\arraystretch}{1.25}
  \begin{tabular}{@{}p{0.30\linewidth}p{0.28\linewidth}p{0.30\linewidth}@{}}
    \toprule
    \textbf{Gap addressed} & \textbf{Approach} & \textbf{Reported gain} \\
    \midrule
    Model fidelity $\rightarrow$ tracking error
      & Complete Lagrangian + linearized motor delay \cite{li2025mecanumdob}
      & DTW $0.2662$~m vs.\ $0.3198$~m ($\sim$17\%) \\
    Disturbance rejection (load, noise)
      & IMC + Disturbance Observer \cite{li2025mecanumdob}
      & 97.6\% / 98.3\% faster settling vs.\ IMC / PID \\
    Inclined surfaces (slip, gravity)
      & ADRC \cite{hernandez2025inclined}
      & 49--60\% better tracking vs.\ PD (experimental) \\
    Unknown CoG offset + slope
      & MRAC, Lyapunov adaptation \cite{chaichumporn2025adaptive}
      & Real-time compensation; stability proof \\
    Disturbance + model uncertainty
      & AVPSMO feedforward in MPC \cite{avpsmo2024}
      & Lyapunov-proven suppression \\
    \bottomrule
  \end{tabular}
  \caption{Higher-fidelity model-based control on Mecanum platforms. All figures
  are the cited authors' own results against their own baselines.}
  \label{tab:control}
\end{table}

\subsection{Deployed control standard versus the research frontier}
\label{sec:related-deployed}

There is a wide gap between what industry ships and what the literature proposes.
Commercial Mecanum platforms are driven, in practice, by kinematic inverse-Mecanum
allocation with per-wheel PID beneath a SLAM/navigation stack, and their positioning
accuracy is secured largely through mechanical precision and external referencing
rather than advanced control---the KUKA omniMove specification of $\pm3$~mm
(base $\pm5$~mm, $\pm1$~mm only with optical/mechanical aids) is representative
\cite{kuka_omnimove,kuka_aerospace}, and youBot-class navigation likewise rests on
kinematic odometry \cite{adamov2018youbot}. The research frontier, by contrast, is
dominated by robust and adaptive designs that repeatedly demonstrate the
nominal/kinematic controller to be insufficient under real disturbance: adaptive
integral terminal sliding mode \cite{aitsm2021}, anti-peaking linear ADRC
\cite{antipeaking2024}, sliding-mode-observer MPC \cite{avpsmo2024}, and adaptive
sliding mode with an extended state observer for automotive intralogistics
\cite{bui2025hybrid}. Yet none of these is standard deployed practice, and---crucially
for this work---none \emph{recovers the per-wheel friction parameters online}. This
is the gap the proposed PINN\,+\,Mamba digital twin targets: closing the
fidelity/identifiability distance between deployed kinematic control and the
disturbance/slip/friction reality the frontier has documented.

\subsection{Physics-informed learning for friction and forward dynamics}
\label{sec:related-pinn}

Physics-informed neural networks (PINNs) are an established route to recovering
friction parameters and learning forward dynamics while preserving physical
consistency. Most directly relevant, a PINN framework fuses the LuGre dynamic
friction model with learnable neural components to identify LuGre parameters from
small, noisy datasets at higher fidelity than simplified friction models, validated
on an underactuated nonlinear system \cite{ozmen2026lugre}---the same
identification problem this paper poses for per-wheel $(\mu,\chi)$ recovery. More
broadly, embedding robot dynamics into the learner improves both identification and
prediction: a PINN that encodes joint dynamics inside a recurrent network with
customized Runge--Kutta cells outperforms conventional PINNs and grey-box
state-space identification and matches ground-truth parameters on UR3e hardware
\cite{yang2023pinnjoint}, and enhanced PINNs that additionally model motor dynamics
and external-force coupling report 31.12\%/37.07\% modeling-accuracy gains and
40.31\%/51.79\% joint-tracking-error reductions over prior state-of-the-art
\cite{fitee2025pinn}. For the nonsmooth behavior central to slip, PINNs that embed
nonsmooth multibody formulations capture stick--slip and separation--reattachment
without the vanishingly small time steps of conventional time-stepping
\cite{li2023nonsmooth}, and Kolmogorov--Arnold-network methods recover interpretable
closed-form static-friction laws ($R^2 > 0.95$) with no predefined functional form
\cite{wang2025kan}. Notably, rigid-body-prior networks alone are insufficient:
standard Deep Lagrangian Networks fail to capture nonlinear friction and
flexibility, motivating augmentations such as DeLaN-FFNN that add a learnable
non-conservative component \cite{li2025compliant}---direct justification for pairing
a physics prior with a learnable friction term rather than relying on either alone.

\subsection{Selective state-space models as a sequence backbone}
\label{sec:related-ssm}

Trajectory data for a richly-excited Mecanum platform are long, and the friction
state (e.g.\ LuGre bristle deflection) is history-dependent, motivating a
sequence model that scales. Mamba introduces selective state-space models with
input-dependent state transitions, achieving linear-time complexity in sequence
length and demonstrated improvements on real data up to million-length sequences
\cite{gu2023mamba}, providing an efficient long-horizon backbone for the PINN +
Mamba digital twin proposed here.

\subsection{The gap this paper addresses}
\label{sec:related-gap}

Two observations define the opening. First, the strongest PINN friction-ID and
forward-dynamics evidence
\cite{yang2023pinnjoint,fitee2025pinn,wang2025kan,li2025compliant} is drawn from
fixed-base industrial manipulators and cobot joints, not from Mecanum
omnidirectional mobile robots; the single most on-point friction-ID result
\cite{ozmen2026lugre} is validated on one underactuated testbed. Their transfer to
\emph{per-wheel} friction $(\mu,\chi)$ recovery on a four-Mecanum platform is thus
established only by analogy. Second, to our knowledge no prior work applies a
selective state-space (Mamba) backbone to Mecanum or mobile-robot forward-dynamics
learning. This paper connects these threads: a physics-informed, Mamba-backed
digital twin that learns forward dynamics and recovers interpretable per-wheel
friction parameters for a KUKA youBot Mecanum platform, targeting exactly the
slip/friction/CoG chokepoint that Section~\ref{sec:related-control} shows to be the
binding constraint on Mecanum control fidelity.

% ===== SECTION CONTENT ENDS ===========================================

\appendix

\section{Certifiability Implications of Interpretable, Physics-Informed Control}
\label{app:certifiability}

Deploying an AMR in close proximity to humans is gated less by nominal tracking
accuracy than by \emph{functional-safety certification}. Under ISO~3691-4:2023,
driverless industrial trucks realize their protective functions---personnel
detection, protective stop, and speed-and-separation monitoring---through the
methodology of ISO~13849-1, whose Performance Levels (PL~a--e) map onto the Safety
Integrity Levels of IEC~61508 \cite{iso3691_4,iso13849,iec61508}. Operation in a
shared human workspace typically requires the relevant safety function to reach
\textbf{PL~d (SIL~2)} in the forward direction (Table~\ref{tab:pl-sil}). The
decisive obstacle for learning-based control is \emph{systematic capability}: these
standards demand evidence that the control software is verifiable, output-bounded,
structurally testable, and free of systematic faults. A black-box neural controller
cannot furnish this evidence---its outputs cannot be exhaustively bounded, its
structure cannot be covered by test, and its behavior cannot be argued
deterministic---which is why systematic reviews find neural controllers
``confined to the lowest levels of safety integrity''
\cite{zhang2020nnverif,certifyml2022}.

\begin{table}[t]
  \centering
  \small
  \renewcommand{\arraystretch}{1.25}
  \begin{tabular}{@{}p{0.13\linewidth}p{0.24\linewidth}p{0.08\linewidth}p{0.42\linewidth}@{}}
    \toprule
    \textbf{Level} & \textbf{$\mathrm{PFH_d}$ (dangerous failures/h)} & \textbf{$\approx$SIL}
      & \textbf{What it permits for an AMR near people} \\
    \midrule
    none / PL~a & $\ge 10^{-5}$ & --- &
      Not safety-creditable; the function may run only \emph{behind} an independent
      certified safety layer, or in segregated/caged space with no humans present. \\
    PL~b & $3{\times}10^{-6}$--$10^{-5}$ & SIL~1 &
      Low-consequence safety functions (warnings, non-critical interlocks). \\
    PL~c & $10^{-6}$--$3{\times}10^{-6}$ & SIL~1 &
      Moderate functions, e.g.\ reverse / low-speed travel where risk is lower. \\
    PL~d & $10^{-7}$--$10^{-6}$ & SIL~2 &
      \emph{Workhorse level for shared workspaces}: personnel detection, protective
      stop, speed-and-separation monitoring, forward travel among humans. \\
    PL~e & $10^{-8}$--$10^{-7}$ & SIL~3 &
      Highest-hazard cases (high mass/speed, severe-injury potential); rarely
      required for standard AMRs. \\
    \bottomrule
  \end{tabular}
  \caption{ISO~13849-1 Performance Levels, their target dangerous-failure rates and
  approximate IEC~61508 SIL mapping, and what each permits for a Mecanum AMR
  operating near humans \cite{iso13849,iec61508,iso3691_4}.}
  \label{tab:pl-sil}
\end{table}

This is precisely where the physics-informed structure of the proposed twin changes
the achievable certification \emph{status}---though not, by itself, the Performance
Level. A model that retains a known friction structure (LuGre) and rigid-body
constraints, identifies interpretable per-wheel parameters $(\mu,\chi)$ with physical
bounds, and admits provable output bounds and reachability supplies exactly the
systematic-capability and verification evidence that a black box cannot
\cite{ozmen2026lugre,wang2025kan}. The effect is a change of role rather than an
automatic level jump: the learned component moves from \emph{non-creditable}
(effectively SIL~0, requiring it to be fully fenced off or overridden) to
\emph{verifiable and creditable} within the safety argument (plausibly SIL~1--2 for
its role). Reaching PL~d/SIL~2 still requires the full ISO~13849 package---Category-3
architecture, diagnostic coverage, and hardware fault tolerance---so interpretability
unlocks the software/systematic half of the case, not the whole of it. In practical
terms it yields two deployment gains: (i) a verifiable, bounded forward-dynamics
model lets the performance channel operate less conservatively (higher speed, tighter
separation) at the \emph{same} certified PL, reducing reliance on external metrology
and nuisance stops; and (ii) the physics-grounded twin can serve as a certifiable
runtime monitor---its divergence from measured motion is a diagnosable, bounded
trip condition---and as the reachable-set oracle a safety assessor requires. We
emphasize that certification is ultimately an assessor's safety-case judgment;
interpretability is an \emph{enabler of the required evidence}, not a certificate in
itself \cite{zhang2020nnverif,certifyml2022}.

\bibliographystyle{unsrtnat}   % ASME: replace with asmems4 (or your template's .bst)
\bibliography{mecanum_pinn_refs}

\end{document}
